\chapter{Marco teórico y estado de la cuestión}
\label{cap:marcoteorico}

Este capítulo revisa el estado del arte en los campos que confluyen en este trabajo. En primer lugar, se describe la disponibilidad de recursos digitales de las lenguas del mundo y sus consecuencias para el procesamiento del lenguaje natural. A continuación, se introduce el concepto de razonamiento metalingüístico y su relación con los problemas de olimpiadas lingüísticas. Finalmente, se revisan los principales conjuntos de evaluación y estrategias de prompting utilizados para estudiar este tipo de razonamiento en modelos de lenguaje de gran escala.

% ─────────────────────────────────────────────────────────────
\section{Lenguas de baja densidad de recursos}
\label{sec:lowresource}
% ─────────────────────────────────────────────────────────────

Se estima que existen entre 6.000 y 7.000 lenguas vivas en el mundo \citep{eberhard2020ethnologue}. Sin embargo, los lenguajes que todos conocemos o nos puedan sonar no tienen la misma presencia digital que otros de los que puede haber escasa información. Mientras que unas lenguas disponen de grandes cantidades de texto, corpus paralelos, diccionarios, herramientas de análisis y sistemas de traducción relativamente desarrollados, otras, apenas cuentan con material escrito accesible o con datos digitalizados suficientes para entrenar modelos de procesamiento del lenguaje natural.

En este contexto se habla habitualmente de lenguas de baja densidad de recursos (\textit{low-resource languages}), y es que la disponibilidad de recursos puede variar mucho de una lengua a otra. Una lengua puede tener textos sin anotar, pero carecer de corpus paralelos, mientras que otra puede disponer de una gramática o de un diccionario elaborado por lingüistas, pero no de datos en formato útil para entrenar modelos y otras apenas tienen presencia escrita o digital. En todos los casos, el problema de fondo es el mismo: no existe suficiente material lingüístico disponible para aplicar directamente los métodos habituales basados en grandes volúmenes de datos.

\citet{joshi2020state} cuantifican esta desigualdad mostrando que más del 88\% de las lenguas del mundo apenas tienen representación en los conjuntos de datos utilizados para entrenar los modelos actuales. Esta falta de representación afecta de forma directa a tareas como la traducción automática, el reconocimiento de voz, el análisis morfosintáctico o la generación de texto. Los modelos actuales tienden a funcionar mejor en lenguas con mucha presencia digital, mientras que su rendimiento disminuye cuando se enfrentan a lenguas con pocos datos disponibles o con estructuras alejadas de las lenguas dominantes en el preentrenamiento. No sería lo mismo pedirle a un modelo que traduzca idiomas con gran cantidad de información digital como el inglés o el chino, que mostrarle una lengua sinotibetana únicamente hablada por pequeños grupos tribales como es el hakhun. Por ello, el problema no es solo técnico. La falta de herramientas para estas lenguas también afecta a su documentación, a su uso en entornos digitales y a la preservación del patrimonio lingüístico y cultural de sus comunidades hablantes. Por este motivo, el trabajo con lenguas de baja densidad de recursos es una línea importante dentro del campo de procesamiento de lenguaje natural actual.

La solución viable para abordar este problema es el aprendizaje en contexto (\textit{in-context learning}). En lugar de reentrenar un modelo con grandes cantidades de datos, se le proporcionan unos pocos ejemplos dentro del propio prompt (petición al modelo) y se observa si es capaz de deducir patrones a partir de ellos. Esta idea encaja especialmente bien con los problemas de olimpiadas lingüísticas, donde el objetivo no es memorizar una lengua, sino inferir reglas a partir de un conjunto reducido de ejemplos.

% ─────────────────────────────────────────────────────────────
\section{Razonamiento metalingüístico}
\label{sec:metalinguistico}
% ─────────────────────────────────────────────────────────────

El razonamiento metalingüístico consiste en observar una lengua como objeto de análisis. No se trata únicamente de usar el lenguaje para comunicarse, sino de fijarse en cómo funciona. Por ejemplo, ver qué elementos se repiten, qué partes de una palabra cambian, cómo se expresan el tiempo, la persona o el número, o qué orden siguen los constituyentes dentro de una frase.

Este tipo de razonamiento aparece de forma clara en los problemas de la Olimpiada Internacional de Lingüística (IOL), una competición internacional dirigida a estudiantes de secundaria en la que se plantean problemas de análisis lingüístico sobre lenguas normalmente desconocidas para los participantes \citep{derzhanski2010linguistics,iol2026problems}. En ellos, el participante recibe un conjunto reducido de frases en una lengua que no conoce, junto con sus traducciones. A partir de esa información debe deducir vocabulario, identificar patrones gramaticales y aplicar las reglas inferidas a nuevas frases. Estos problemas, como ya se ha mencionado, son autocontenidos: toda la información necesaria para resolverlos está en el propio enunciado. Por ello, este formato resulta especialmente útil para evaluar modelos de lenguaje de gran escala, ya que obliga a trabajar con pocos datos y a razonar sobre regularidades lingüísticas. La tarea no consiste únicamente en producir una traducción que parezca plausible, sino en mantener hipótesis coherentes y aplicarlas de forma sistemática. Por ejemplo, si el modelo identifica que una partícula marca pasado, debería aplicar esa hipótesis a todas las frases compatibles, y no solo a una de ellas. Esto se aplicaría tanto a los tiempos verbales como a otros rasgos gramaticales, como el género, el número, la conjugación verbal o incluso la diferencia entre una pregunta y una oración.

Aun así, es importante distinguir entre el resultado que observamos y el proceso interno del modelo. Que un modelo acierte una traducción no implica necesariamente que haya razonado de la misma forma que una persona. Puede haber aprovechado asociaciones parciales, patrones aprendidos durante el preentrenamiento o simplemente haber generado una respuesta plausible que se asemeja a la solución real. Por eso, en este proyecto nos interesa analizar también el diccionario, las reglas y las explicaciones que el modelo produce durante la resolución.

% ─────────────────────────────────────────────────────────────
\section{Benchmarks de olimpiadas lingüísticas para LLMs}
\label{sec:benchmarks}
% ─────────────────────────────────────────────────────────────

En los últimos años han aparecido varios conjuntos de evaluación que utilizan problemas de olimpiadas lingüísticas para estudiar el razonamiento lingüístico de los modelos de lenguaje. Todos ellos parten de una idea parecida: proporcionar al modelo unos pocos ejemplos en una lengua poco conocida y comprobar si es capaz de extraer regularidades, formular hipótesis y aplicarlas a nuevas expresiones. Este tipo de evaluación es nos sirve para estudiar el aprendizaje en contexto: a diferencia de otras tareas de traducción o comprensión lingüística, aquí no interesa medir únicamente cuánto sabe previamente el modelo sobre una lengua concreta, sino hasta qué punto puede razonar a partir de la información incluida en el propio enunciado. Por este motivo, estos benchmarks son una referencia directa para el diseño experimental de este trabajo.

A continuación se revisan los principales conjuntos de datos relacionados con esta línea de investigación.

\subsection{PuzzLing Machines}

El trabajo de \citet{sahin2020puzzling} es uno de los primeros intentos de evaluar sistemas automáticos mediante problemas de olimpiadas lingüísticas. Los autores recopilan un conjunto cercano a cien problemas de tipo piedra Rosetta, procedentes de olimpiadas lingüísticas, que cubren 81 lenguas pertenecientes a 39 familias distintas. A partir de este conjunto, analizan hasta qué punto los modelos de traducción automática son capaces de resolver tareas que requieren aprender a partir de muy pocos ejemplos, por lo que diferencia de una tarea de traducción automática convencional, estos problemas no proporcionan un corpus amplio, sino un conjunto mínimo de expresiones paralelas a partir del cual el sistema debe inferir léxico, reglas gramaticales y fenómenos morfológicos o fonológicos. Los autores evalúan aproximaciones de traducción automática previas a la aparición de los LLMs conversacionales actuales. En concreto, comparan métodos estadísticos simples con modelos neuronales de traducción de tipo secuencia a secuencia y arquitecturas profundas de traducción automática. Estos sistemas se entrenan o ajustan sobre los pocos pares disponibles en cada puzzle, de forma que el experimento mide si son capaces de generalizar desde un conjunto de entrenamiento extremadamente pequeño. Por este motivo, el trabajo resulta especialmente relevante como antecedente: no evalúa modelos como podrían ser GPT o Claude, sino técnicas que hoy pueden considerarse en gran parte obsoletas frente a los LLMs actuales.

Los resultados muestran que tanto los métodos estadísticos como los modelos neuronales evaluados tienen un rendimiento insuficiente en este tipo de problemas. De los experimentos realizados, se deduce que los modelos de traducción basados en patrones superficiales y grandes cantidades de datos no son adecuados para puzzles que requieren razonamiento metalingüístico, deducción y revisión de hipótesis a partir de unos pocos ejemplos. Esta conclusión es importante para este trabajo porque justifica el paso de la traducción automática convencional hacia estrategias de prompting más explícitas, donde el modelo debe construir un diccionario y reglas gramaticales antes de traducir.

\subsection{modeLing}

\citet{chi2024modeling} presentan \textit{modeLing}, un conjunto de problemas construido manualmente y diseñado específicamente para evaluar razonamiento lingüístico reduciendo el riesgo de contaminación de datos. A diferencia de otros benchmarks basados directamente en problemas ya publicados de la IOL, modeLing utiliza problemas creados para el propio trabajo, lo que disminuye la posibilidad de que los modelos hayan visto los enunciados durante el preentrenamiento.

El conjunto contiene 48 problemas y 272 preguntas, basados en 19 lenguas con presencia digital muy reducida. Los ejercicios cubren cuatro tipos principales de razonamiento, que son el orden entre nombre y adjetivo, orden de constituyentes, posesión y semántica. Además, los autores realizan una prueba sin ejemplos de referencia para detectar posible contaminación de datos. En la evaluación se prueban seis modelos de la familia GPT: Ada, Babbage, Curie, Davinci, GPT-3.5 y GPT-4. Los cuatro primeros pertenecen a la familia GPT-3; tomando como referencia los tamaños descritos por \citet{brown2020language}, Ada, Babbage, Curie y Davinci corresponden aproximadamente a modelos de 350M, 1.3B, 6.7B y 175B parámetros, respectivamente. En cambio, el tamaño exacto de GPT-3.5 y GPT-4 no es público, por lo que el trabajo los trata como modelos cerrados de mayor capacidad. Cabe mencionar que esta aclaración es importante porque varios de los modelos evaluados, así como Ada, Babbage y Curie, han quedado obsoletos frente a los modelos actuales. Los autores comparan cuatro configuraciones de prompting: un prompt mínimo, un prompt ajustado manualmente por una persona medallista de la Olimpiada Internacional de Lingüística, un \textit{chain-of-thought} básico y un \textit{full chain-of-thought} con un ejemplo de razonamiento paso a paso. Los resultados muestran una diferencia clara entre modelos pequeños y grandes. Ada, Babbage y Curie quedan prácticamente en cero en todas las configuraciones, con valores máximos de 0.011, 0.018 y 0.022 de exactitud, respectivamente. En cambio, Davinci alcanza hasta 0.514 con \textit{full chain-of-thought}, GPT-3.5 obtiene valores alrededor de 0.40 y GPT-4 alcanza el mejor resultado, con 0.607 de exactitud en la configuración \textit{full chain-of-thought}. 

Estos resultados son relevantes para este trabajo por dos motivos. En primer lugar, muestran que los modelos grandes sí pueden resolver parte de los puzzles mediante razonamiento con pocos ejemplos, aunque el rendimiento sigue lejos de ser perfecto. En segundo lugar, indican que cambiar el tipo de prompt no siempre produce mejoras grandes: en \textit{modeLing}, las diferencias entre prompt mínimo, prompt ajustado y razonamiento paso a paso son menores que las diferencias entre tamaños de modelo. Esto refuerza la necesidad de analizar conjuntamente el efecto de la metodología y el modelo utilizado.

\subsection{Linguini}

\citet{sanchez2024linguini} presentan \textsc{Linguini}, el benchmark utilizado como dataset base de mi trabajo. Está formado por 894 preguntas agrupadas en 160 problemas extraídos del corpus de la IOL, y cubre 75 lenguas de baja densidad de recursos. Los problemas se organizan en tres categorías principales:

\begin{enumerate}
	\item \textbf{Transducción de secuencias} (\textit{sequence transduction}): el modelo debe transformar una secuencia en otra, por ejemplo traduciendo una frase o relacionando expresiones equivalentes.
	\item \textbf{Relleno de huecos} (\textit{fill-in-blanks}): el modelo debe completar formas lingüísticas aplicando reglas morfológicas o fonológicas deducidas a partir de los ejemplos.
	\item \textbf{Transliteración numérica} (\textit{number transliteration}): el modelo debe convertir expresiones numéricas entre su representación textual y su representación mediante cifras.
\end{enumerate}

Los resultados publicados muestran que estos problemas siguen siendo difíciles para los modelos actuales. Ninguno de los modelos evaluados supera el 25\% de exactitud. El mejor modelo cerrado, Claude 3 Opus, alcanza un 24,05\% de exactitud, mientras que el mejor modelo abierto, Llama 3 70B, obtiene tan solo un 8,84\%. Como aclaración, el término de exactitud (Exact Match) se refiere al porcentaje de respuestas que coinciden exactamente con la referencia esperada, por ello una respuesta parcialmente correcta puede recibir una puntuación baja si no coincide de forma exacta con la solución del benchmark. Como métrica complementaria, los autores también reportan chrF, que permite medir similitudes parciales a nivel de caracteres. Los valores resultantes quedan por debajo del rendimiento humano reportado en competiciones de la IOL, donde estudiantes de secundaria o preuniversitarios suelen superar los 82 puntos en los mejores resultados anuales \citep{sanchez2024linguini,iol2026problems}. Esta comparación debe interpretarse teniendo en cuenta el perfil de los participantes, y es que la IOL está dirigida a estudiantes preuniversitarios, habitualmente de entre 15 y 20 años.

El experimento más relevante de Linguini es el de la prueba de la ablación sin contexto. Al eliminar las frases de entrenamiento del prompt, el rendimiento de los modelos cae de forma clara. Por ejemplo, Claude 3 Opus pasa de 24,05 a 1,23, GPT-4o de 14,65 a 1,45 y Llama 3 70B de 8,17 a 1,67. Aunque parezca obvio, esto es una demostración clara de que la información proporcionada en el contexto tiene un papel crucial en la resolución. Aun así, no es posible descartar por completo la influencia del conocimiento adquirido durante el preentrenamiento, porque los datos exactos usados para entrenar muchos de estos modelos no son públicos. Por tanto, una caída fuerte sin contexto es un indicio de baja contaminación, pero no una prueba absoluta de que el modelo no haya visto nunca información relacionada con la lengua o con el problema.

\subsection{LingOly}

\citet{bean2024lingoly} presentan \textit{LingOly}, un benchmark de mayor escala que los conjuntos revisados anteriormente, especialmente en comparación con \textit{PuzzLing Machines}, \textit{modeLing} y \textsc{Linguini}, publicado en NeurIPS 2024. El conjunto cubre más de 90 lenguas, incluidas algunas lenguas extintas, y contiene 1.133 problemas organizados en seis categorías lingüísticas: fonología, morfología, sintaxis, semántica, sistemas numéricos y problemas compuestos. Además, los problemas se clasifican en cinco niveles de dificultad humana.

Esta organización permite analizar con más detalle qué tipos de fenómenos resultan más difíciles para los modelos. Por ejemplo, el benchmark incluye problemas relacionados con cambios fonológicos, formación de palabras, orden de constituyentes, concordancia, posesión, sistemas de numeración y combinaciones de varios fenómenos dentro de un mismo ejercicio. Los autores evalúan once modelos de lenguaje y comparan sus resultados con un baseline sin contexto. Los resultados muestran que los modelos tienen especiales dificultades en los problemas de mayor nivel. Incluso el mejor sistema evaluado alcanza únicamente un 38,7\,\% de exactitud en los ejercicios más complejos, aunque mejora en 24,7 puntos respecto a la versión sin contexto.

El trabajo también observa una tendencia general: los modelos propietarios de mayor tamaño suelen superar a los modelos abiertos, y el rendimiento tiende a ser mayor en lenguas con más presencia digital.

\subsection{Comparación entre los principales benchmarks}

Los cuatro conjuntos revisados tienen objetivos complementarios. \textit{PuzzLing Machines} plantea uno de los primeros marcos para usar problemas de tipo piedra Rosetta como prueba de aprendizaje con pocos datos. \textit{modeLing} se centra en reducir el riesgo de contaminación mediante problemas creados expresamente para la evaluación. \textit{Linguini} amplía la escala del análisis y proporciona el marco experimental utilizado en este trabajo. Por último, \textit{LingOly} incorpora una mayor variedad de formatos, lenguas y niveles de dificultad.

La Tabla~\ref{tab:comparacion_benchmarks} resume las características principales de estos benchmarks. En el caso de \textit{PuzzLing Machines}, el artículo describe el conjunto como cercano a 100 puzzles, mientras que el conteo del dataset utilizado corresponde a 96 problemas y 924 preguntas.

\begin{table}[ht]
	\centering
	\caption{Comparación de los principales benchmarks de olimpiadas lingüísticas para LLMs.}
	\label{tab:comparacion_benchmarks}
	\begin{tabularx}{\textwidth}{l c c X}
		\toprule
		\textbf{Benchmark} & \textbf{Año} & \textbf{Tamaño} & \textbf{Rasgo principal} \\
		\midrule
		\textit{PuzzLing Machines} 
		& 2020 
		& 96 problemas, 924 preguntas
		& Antecedente centrado en problemas de tipo piedra Rosetta. \\
		
		\textit{modeLing} 
		& 2024 
		& 48 problemas, 272 preguntas 
		& Problemas originales diseñados para reducir el riesgo de contaminación de los datos de entrenamiento. \\
		
		\textit{Linguini} 
		& 2024 
		& 160 problemas, 894 preguntas 
		& Benchmark de referencia utilizado en los experimentos de este trabajo. \\
		
		\textit{LingOly} 
		& 2024 
		& 1.133 preguntas individuales
		& Conjunto de mayor tamaño, con distintos formatos y niveles de dificultad. \\
		\bottomrule
	\end{tabularx}
\end{table}

En conjunto, estos trabajos muestran que los problemas de olimpiadas lingüísticas son una herramienta útil para estudiar la capacidad de los modelos de lenguaje para generalizar a partir de pocos ejemplos. También dejan claro que el rendimiento sigue siendo limitado y que la forma de presentar la información al modelo puede influir de manera importante en los resultados. Esta última idea es especialmente relevante para este trabajo, ya que la aportación principal no consiste en proponer un nuevo benchmark, sino en comparar distintas formas de guiar el razonamiento del modelo sobre los mismos problemas.

% ─────────────────────────────────────────────────────────────
\section{Estrategias de prompting para razonamiento lingüístico}
\label{sec:prompting}
% ─────────────────────────────────────────────────────────────

Ante los resultados todavía limitados del prompting directo en tareas de razonamiento lingüístico, varios trabajos han explorado estrategias más elaboradas para guiar al modelo en su razonamiento y en sus conclusiones. En los benchmarks revisados anteriormente, estas mejoras se han planteado principalmente mediante cambios en la forma de presentar el problema: añadir ejemplos de razonamiento, organizar los pasos de análisis, introducir instrucciones más explícitas o pedir al modelo que revise sus hipótesis antes de responder \citep{chi2024modeling,zhu2025evaluating,madaan2023selfrefine}. Esta sección no se limita únicamente a trabajos sobre olimpiadas lingüísticas, sino que recoge también estrategias generales de prompting que resultan relevantes para este tipo de tareas. En todos los casos, se trata de métodos que no modifican los parámetros del modelo, ya que el objetivo es evaluar qué puede hacer un modelo ya entrenado a partir de la información incluida en el prompt, sin recurrir a fine-tuning ni a entrenamiento adicional. Esta decisión es especialmente importante en lenguas de baja densidad de recursos, donde no suele haber datos suficientes para ajustar un modelo de forma fiable.

\subsection{Prompting analógico}

\citet{ramji2024inductive} proponen una estrategia basada en el razonamiento analógico. En lugar de depender únicamente de los ejemplos incluidos en el problema, el modelo genera o recupera ejemplos auxiliares relacionados con fenómenos lingüísticos similares. Dicho de otra manera, el modelo trata de aprovechar el conocimiento general adquirido durante el preentrenamiento. Por ejemplo, si una lengua desconocida presenta una una estructura sintáctica poco habitual, el modelo puede intentar relacionarlo con estructuras semejantes presentes en otras lenguas que conoce mejor.

Esta estrategia puede facilitar la identificación de determinados patrones, pero también introduce una limitación importante, además de desviarse de nuestro objetivo principal, y es que el razonamiento deja de depender exclusivamente de la información proporcionada en el enunciado. Otro problema es que el modelo, al igual que saca conclusiones verídicas, también puede apoyarse en analogías incorrectas o trasladar reglas de una lengua a otra sin suficiente evidencia.

En mi investigación no se utiliza prompting analógico, ya que uno de los objetivos es evaluar hasta qué punto el modelo puede resolver los problemas a partir de los ejemplos disponibles, sin recibir pistas sobre la lengua ni sobre su familia lingüística. No obstante, esta aproximación podría resultar interesante como posible línea de ampliación futura.

\subsection{Gramática en contexto}

\citet{tanzer2024benchmark} estudian otra aproximación: proporcionar al modelo una gramática descriptiva de la lengua como parte del contexto. Este enfoque, denominado \textit{one-book prompting}, permite trabajar con lenguas para las que existe documentación lingüística, aunque no se disponga de grandes corpus textuales. A diferencia de los problemas de olimpiadas lingüísticas, el modelo no debe deducir todas las reglas exclusivamente a partir de ejemplos. En su lugar, puede consultar una descripción explícita de la lengua y utilizarla para traducir nuevas frases.

Los resultados muestran que incluir una gramática puede mejorar el rendimiento del modelo, especialmente cuando se combina con ejemplos. Resulta que muchas lenguas con pocos recursos digitales sí disponen de gramáticas, diccionarios o estudios elaborados por especialistas, por lo que esta estrategia parece prometedora para aplicaciones reales de traducción y documentación lingüística. Sin embargo, este enfoque responde a un objetivo diferente del que se pretende abordar. Aquí el objetivo está en evaluar la capacidad de deducción a partir de un conjunto reducido de pares de traducción, sin incorporar documentación externa sobre la lengua.

Sin embargo, este enfoque también ha recibido críticas. \citet{aycock2025grammarbook} señalan que la mejora observada al añadir una gramática no siempre se debe a que el modelo esté usando realmente las explicaciones gramaticales. En muchos casos, la mejora parece venir sobre todo de los ejemplos de traducción que acompañan a la gramática. Es decir, el modelo aprovecha mejor los pares de frases ya traducidas que las reglas lingüísticas descritas en el documento. Los autores indican que una gramática puede ser útil para tareas como decidir si una frase está bien formada o predecir glosas, pero no encuentran pruebas claras de que los LLMs sepan usar una gramática completa para traducir mejor. Esta crítica es importante para este trabajo porque obliga a distinguir bien qué tipo de ayuda recibe el modelo: ejemplos paralelos, reglas explícitas o instrucciones para razonar paso a paso.

\subsection{Aprendizaje en contexto con ejemplos}

El aprendizaje en contexto (\textit{in-context learning}, ICL) es la capacidad de un modelo de lenguaje para realizar una tarea nueva a partir de la información incluida directamente en el prompt, sin necesidad de actualizar sus parámetros. Esta capacidad se acerca más a lo que pretendemos, ya que en los problemas de olimpiadas lingüísticas, la propiedad in-context resulta fundamental. El modelo recibe un conjunto reducido de pares formados por una expresión en una lengua desconocida y su traducción. A partir de estos ejemplos debe deducir vocabulario, identificar regularidades morfológicas y aplicar las reglas extraídas a nuevas frases.

Es importante diferenciar este mecanismo del \textit{few-shot prompting}. En una configuración \textit{few-shot}, el modelo recibe ejemplos adicionales que muestran cómo debe resolverse la tarea o cómo debe estructurarse la respuesta. Estos ejemplos no forman parte necesariamente del problema lingüístico concreto, sino que actúan como demostraciones previas. Siempre que se trate de dar información adicional como apoyo al modelo, desviamos el objetivo principal de comprobar la capacidad de razonamiento metalingüístico de los modelos de lengua. 

\citet{sanchez2024linguini} comparan distintas configuraciones con un número variable de demostraciones y observan resultados extraídos de la respuesta de los modelos. Algunos de ellos mejoran al recibir ejemplos adicionales, mientras que otros empeoran. La explicación más plausible hasta el momento es que las demostraciones introducen información útil sobre el formato de la tarea, pero también pueden añadir ruido cuando proceden de lenguas o estructuras poco relacionadas con el problema que se desea resolver. Esta observación es bastante interesante para la presente investigación: proporcionar más información no garantiza por sí solo una mejora del rendimiento. Es crucial estudiar cómo se presenta esa información, en qué orden se incorpora y si el modelo es capaz de mantener hipótesis coherentes a lo largo de la resolución.

\subsection{Razonamiento paso a paso}

\citet{zhu2025evaluating} proponen una estrategia de descomposición paso a paso inspirada en el \textit{chain-of-thought prompting} \citep{wei2022chain}. La idea general consiste en dividir un problema complejo en pasos intermedios más sencillos de manejar, en lugar de pedir directamente la respuesta final. En el contexto de las olimpiadas lingüísticas, este razonamiento puede consistir en identificar primero el vocabulario básico, después analizar los patrones morfológicos, formular reglas gramaticales y, finalmente, aplicar esas reglas a las frases de prueba.

En el estudio de la descomposición progresiva, en lugar de presentar todos los ejemplos de entrenamiento de golpe, los autores organizan la información en etapas lingüísticas (concretamente, en 4 pasos): primero se presenta el léxico y el orden de palabras básico, después los patrones fonológicos, luego la morfosintaxis y finalmente la sintaxis compleja. En cada etapa, el modelo actualiza su comprensión de la lengua antes de pasar a la siguiente. Este enfoque obtiene mejoras significativas respecto al prompting directo en lenguas con morfología verbal compleja. Aquí comprobaremos dicha estrategia, ya que resulta especialmente observar que la forma de ordenar y estructurar los ejemplos puede influir en el rendimiento del modelo, incluso cuando la información lingüística disponible es la misma.

\subsection{Autocorrección y revisión de respuestas}

Otra línea de trabajo relacionada es la autocorrección iterativa. El paradigma \textit{Self-Refine}, propuesto por \citet{madaan2023selfrefine}, plantea que un modelo puede generar una primera respuesta, producir una crítica sobre ella y después elaborar una versión revisada a partir de esa crítica. Aplicado a tareas de razonamiento lingüístico, este planteamiento permite estudiar si el modelo es capaz de detectar contradicciones internas en su propia respuesta. La idea es que si se formula una regla gramatical durante el análisis, pero después genera una traducción que no respeta esa regla, una fase de revisión debería permitir identificar el problema y corregirlo.

Este enfoque puede tener una limitación importante, que es la retroalimentación procede del propio modelo. Si la primera interpretación contiene una regla equivocada, el modelo puede utilizar esa misma regla como criterio de revisión y reforzar el error en lugar de corregirlo, aunque se le indique explícitamente que esta primera hipótesis puede contener errores. Por tanto, la autocorrección puede ser útil en algunos casos, pero no garantiza necesariamente una mejora del resultado final.

La estrategia de la autocorrección también se incorpora en esta investigación, se diseña prompt que pida una segunda fase de revisión, en la que el modelo debe comprobar la coherencia entre sus traducciones, su diccionario y las reglas que ha formulado. Como veremos más adelante, los resultados no siempre mejoran respecto al baseline, lo que permite analizar también las limitaciones de este tipo de revisión.

% ─────────────────────────────────────────────────────────────
\section{Análisis de las dificultades de los LLMs en puzzles lingüísticos}
\label{sec:dificultades}
% ─────────────────────────────────────────────────────────────

Además de medir cuántas respuestas aciertan los modelos, varios trabajos recientes han intentado analizar por qué fallan en los problemas de olimpiadas lingüísticas. Esta cuestión es importante porque dos modelos pueden obtener una puntuación parecida y, sin embargo, cometer errores de base muy diferentes. Mientras que uno puede fallar por no identificar vocabulario básico, otro por no mantener una regla morfológica de forma consistente y otro por producir una traducción gramaticalmente plausible pero incompatible con los ejemplos del enunciado.

\citet{choudhary2025unveiling} realizan uno de los análisis más completos en esta dirección. En su trabajo analizan el rendimiento de modelos de lenguaje sobre 629 problemas en 41 lenguas de baja densidad de recursos, anotando cada puzzle con 50 características lingüísticas. Entre los modelos evaluados se incluyen sistemas cerrados y abiertos, como GPT-4o, Claude 3.5 Sonnet, Llama 3.1 8B, Llama 3.1 70B, Gemma 2 9B, Gemma 2 27B, Mistral 7B, Mixtral 8x7B y Qwen 2.5 7B. Las conclusiones que sacan resultan muy reveladoras y apuntan a dos factores principales. Por un lado, deducen que la complejidad morfológica parece ser uno de los elementos que más dificulta la resolución, y es que cuanto mayor es el número de morfemas por palabra, peor tienden a comportarse los modelos. Por otro, los modelos obtienen mejores resultados cuando las características del puzzle se parecen a patrones presentes en inglés, en español o en otras lenguas con alta representación digital. Además, los autores muestran que separar previamente las palabras en morfemas mejora el rendimiento, lo que sugiere que parte del problema está relacionado con la forma en que los modelos tokenizan y representan lenguas poco familiares \citep{choudhary2025unveiling}. Con ello, cabe destacar que no todos los puzzles tienen la misma dificultad, aunque tengan un número parecido de ejemplos para el entrenamiento. Una lengua con marcas de persona, número, caso o tiempo muy fusionadas en la palabra puede ser más difícil que otra con una segmentación más transparente. Por ejemplo, es más sencillo identificar el plural si las palabras que lo indican siempre contienen la terminación `s', o el género si lo que cambia es el artículo anterior a la misma palabra. Por tanto, las diferencias de rendimiento entre lenguas no deben explicarse únicamente por el modelo o por la estrategia de prompting utilizada, sino también por las propiedades lingüísticas del propio puzzle, como también comprobaremos en los resultados de la investigación.

Un trabajo posterior de \citet{majmudar2025llms} amplía la evaluación a modelos de razonamiento y estudia tanto la resolución como la generación de puzzles lingüísticos. El trabajo utiliza problemas de olimpiadas lingüísticas dirigidos a estudiantes de secundaria y evalúa modelos recientes, incluyendo o1 de OpenAI, junto con otros LLMs de alto rendimiento. Los resultados muestran que los modelos de razonamiento mejoran respecto a sistemas anteriores en varias categorías de puzzles, e incluso pueden superar a participantes humanos en algunos tipos de ejercicios. Sin embargo, los autores señalan que siguen apareciendo dificultades claras en problemas relacionados con sistemas de escritura poco familiares y con lenguas poco estudiadas o tipológicamente alejadas de las lenguas más presentes en los datos de entrenamiento \citep{majmudar2025llms}.

Este resultado sugiere que aumentar la capacidad de razonamiento del modelo puede mejorar el rendimiento, pero no elimina por completo las limitaciones asociadas a la distancia tipológica, la morfología compleja o la falta de familiaridad con determinados sistemas lingüísticos. Además, el propio trabajo muestra que resolver puzzles y generarlos no son tareas equivalentes: los modelos pueden acertar algunos problemas, pero tienen más dificultades para crear puzzles nuevos que sean correctos, coherentes y resolubles por una persona.

En conjunto, estos trabajos muestran que los errores de los modelos no son aleatorios, sino que suelen concentrarse en fenómenos concretos: marcas morfológicas pequeñas, distinciones de número o persona, alternancias internas de la palabra, sistemas de escritura no familiares o reglas que deben aplicarse de forma consistente a varias frases. Esta observación explica que, además de calcular métricas automáticas, sea necesario revisar cualitativamente algunos casos para entender qué tipo de razonamiento está realizando el modelo y dónde se rompe.

% ─────────────────────────────────────────────────────────────
\section{Síntesis y posicionamiento del trabajo}
\label{sec:sintesis}
% ─────────────────────────────────────────────────────────────

La revisión anterior permite situar el presente trabajo de fin de grado dentro de una línea de investigación reciente: la evaluación del razonamiento lingüístico de los modelos de lenguaje a partir de problemas de olimpiadas lingüísticas. Los trabajos revisados muestran que los modelos actuales pueden identificar algunos patrones a partir de pocos ejemplos, pero también que su rendimiento sigue estando lejos de ser estable. La dificultad depende de muchos factores: la lengua del puzzle, la complejidad morfológica, la dirección de traducción, el formato del prompt y la capacidad del modelo para mantener hipótesis coherentes.

A partir de este contexto, aquí no se propone un nuevo benchmark, sino que utiliza Linguini como base experimental para comparar distintas formas de presentar la información al modelo. El objetivo principal no es solo medir cuántas traducciones son correctas, sino observar si una estrategia de prompting más estructurada ayuda al modelo a construir un diccionario, formular reglas gramaticales y revisar sus propias hipótesis durante la resolución.

La metodología de este trabajo se parece a los estudios previos en que también evalúa modelos de lenguaje mediante problemas lingüísticos autocontenidos y utiliza ejemplos dentro del prompt para inducir reglas, como ocurre en benchmarks como \textit{PuzzLing Machines}, \textit{modeLing}, \textsc{Linguini} o \textit{LingOly}. También comparte con trabajos sobre chain-of-thought, autocorrección y razonamiento guiado la idea de que no basta con pedir una respuesta final, sino que puede ser útil hacer explícitos pasos intermedios del análisis. Sin embargo, se diferencia de esos trabajos en que no propone nuevos problemas ni se centra únicamente en comparar modelos, sino en comparar varias metodologías de presentación y revisión sobre el mismo conjunto de puzzles. En particular, se estudian estrategias como la presentación incremental acumulativa, la inspiración en una resolución humana y la autocorrección iterativa, manteniendo constante el dataset para observar mejor el efecto del diseño del prompt.

La línea de trabajo general consiste en establecer varias metodologías, diferentes entre sí, y observar como se comportan con puzzles de distintas familias y estructuras. Con estas estrategias, nos situaremos entre la evaluación cuantitativa y el análisis cualitativo. Por un lado, se comparan modelos y metodologías mediante métricas automáticas y, por otro, se revisan casos concretos para entender qué hipótesis formula el modelo, qué errores aparecen y cómo evoluciona su razonamiento a lo largo del proceso.

Cabe destacar que la finalidad no es demostrar que los modelos razonan de la misma manera que una persona, sino analizar si una organización más guiada del proceso de resolución puede producir respuestas más correctas, coherentes y revisables en problemas de olimpiadas lingüísticas.